%======================================================================
\subsection{BW-PDM-12 --- Built 31 Oct: renumber an item, and assemblies numbered as parts}\label{sec:pdm12-built}
%======================================================================

You asked for the battery to be a Part number in Standard Libraries and for the vault to be able to rename an item.
Both are built, as planned in \S\ref{sec:pdm12-plan}. Nothing has been run on the real vaults: that is your step
(below).

\paragraph{What you can do now.}
\begin{itemize}[nosep]
  \item \textbf{Number assemblies as parts.} A box in Vault Settings (Part numbering), stored in the vault's
        \file{numbering.json} as \code{"assemblies\_as\_parts": true}, and
        \code{bw-hooks vault settings --assemblies-as-parts on|off}. Off in a new vault. Import, New Part in Vault,
        project setup, the KiCad import, the connector library and Renumber's suggestion all pick the document type
        through \code{vault\_setup.document\_type\_for\_item}. Drawings keep type 03; the item stays an assembly.
  \item \textbf{Renumber\ldots} in the right-click menu of the Vault panel (and the Vault menu). The dialog
        suggests the next free number, takes a typed one, requires a reason, and shows what will be rewritten, the
        files that link the item and the items in the neighbour vaults that link its file name. It closes the
        affected open documents when you agree and refuses while one has unsaved changes.
  \item \textbf{Command line.} \code{bw-hooks vault renumber PN [--to NEW | --document-type CODE] --reason TEXT
        [--dry-run]}. The dry run needs no FreeCAD. Other vault commands accept an old number and say what it is now.
  \item \textbf{Activity log.} ``renumbered \code{A} $\to$ \code{B}: reason'', with the new file versions in the
        detail. \textbf{Search} finds an item by a former number.
\end{itemize}

\paragraph{How it works.} One new journal operation, \code{renumber} (old, new, reason, and the file versions it
wrote as P/N $\to$ object record). \code{bruceworx.renumber} does the file step first: the item's latest file is
copied, \code{write\_part\_number} writes the new \code{PartNumber} keeping the revision, the link errors the
headless save leaves are cleared, the view data FreeCADCmd drops is restored
(\code{connector\_library.restore\_view\_data}), and the geometry fingerprint is compared (refused if it differs).
It is stored as \file{\textless new\textgreater.FCStd}. Each file that links the item has its link re-pointed
(\code{rewrite\_links}) and is stored as a new version of its current revision. The objects are written before the
single journal line. Replay (\code{pdm.Vault.renumber}, \code{FolderVault}) re-keys the item, versions, locks,
released files, BOM lines in every revision, drawings, configurations, ECR/ECO lists and project tops; the old
number is an alias (\code{vault.resolve}), chained renumbers resolve to the last. Get Revision and link fixing read
old links through the aliases; a read-only work copy under the old name is replaced at the next refresh.

\paragraph{Decided while building.}
\begin{itemize}[nosep]
  \item The journal line carries each new file as \code{\{sha256, size, filename\}}, not the SHA-256 alone, so a
        replay needs no file access. The activity detail shows P/N and the start of the SHA-256.
  \item A released linker's \emph{released} file moves to the new version too (it is the same revision, only the
        link changed); otherwise Released would point at a file that links a number that no longer exists.
  \item The new versions keep the old fingerprint. A drawing's \code{made\_from} is re-keyed to the new number, and
        its entries for rewritten files point at the new versions.
  \item The number is reserved in \file{numbering.json} just before the journal line (a typed number too), so a
        later new number passes it by. A failure before that burns nothing; one after leaves an unused number.
  \item \code{renumber} sits in the contract table as row T16 and in \file{docs/pdm-openapi.yaml}.
  \item The fingerprint reader is a second copy in \code{freecad\_io} (the Vault panel's lives in the GUI module).
\end{itemize}

\paragraph{Every Keystone needs this code.} Like \code{set\_category}, \code{renumber} is a new journal operation:
an older Keystone opening a vault with the line reports it as a problem and goes on with the \emph{old} numbers.
Install the new build on every Mac before using it on a shared vault.

\paragraph{Tests.} \code{test\_renumber.py} (model, every state, refusals, chained aliases; folder vault files and
links, released linker, crash before the line, dry run, KiCad refusal, cross-vault list, activity text, workspace
refresh and Get Revision, the setting in every numbering path, the CLI) and contract row T16; GUI tests in
\code{TestBruceWorxVaultClient} (assembly numbered as a part, renumber with its file, unsaved document refused,
clean one closed, search by old number) and \code{TestBruceWorxVaultWizard} (the box). Tried on a copy of Standard
Libraries: the battery went to the next \code{STD-01} number under FreeCADCmd with its eight links intact.

\paragraph{Not done.} KiCad items are refused. The activity detail for other vaults' links is a list in the dialog
only. Links in the neighbour vaults are not rewritten: they resolve through the aliases and are re-pointed at the
item's next check-in.

\paragraph{Your step.} After the new Keystone is installed, in Standard Libraries:
\begin{verbatim}
bw-hooks vault settings --vault "/Volumes/Extreme SSD/VectorLabs/Vault/Standard Libraries" \
    --assemblies-as-parts on
bw-hooks vault renumber STD-02-01001-001 --vault "/Volumes/Extreme SSD/VectorLabs/Vault/Standard Libraries" \
    --reason "library items are numbered as parts (31 Oct)"
\end{verbatim}
